\section*{Chapter \ref{ch:g6:living-standards} --- Two Hundred Years of Living Standards}

\begin{solution}{exo:g6:living-standards:1}
In 1700, a little above the 1\,000 gridline's ``2'' tick: about 2\,400 dollars.
In 2000, between 10\,000 and 100\,000, a little above the ``3'' tick: about
32\,000 dollars. (The table gives $2\,412$ and $31\,946$.)
\end{solution}

\begin{solution}{exo:g6:living-standards:2}
(a) The Western Offshoots ($2\,513$ dollars). (b) The richest was again the
Western Offshoots ($56\,568$ dollars); the poorest, Sub-Saharan Africa
($3\,437$ dollars).
\end{solution}

\begin{solution}{exo:g6:living-standards:3}
$16\,677 \div 1\,128 \approx 14.8$: almost fifteen times larger.
\end{solution}

\begin{solution}{exo:g6:living-standards:4}
(a) Argentina, by far ($7\,949$ dollars against $998$). (b) At the end of the
1980s: in 1988.
\end{solution}

\begin{solution}{exo:g6:living-standards:5}
(a) The change is $16\,352 - 489 = 15\,863$ dollars, and $15\,863 \div 489 \approx
32.44$: a rise of about $3\,244\,\%$. (b) $16\,352 \div 489 \approx 33.4$ times.
\end{solution}

\begin{solution}{exo:g6:living-standards:6}
(a) $0.22 \times 60 = 13.2$ minutes. (b) $5.37 \div 0.22 \approx 24.4$: about
24 times less work.
\end{solution}

\begin{solution}{exo:g6:living-standards:7}
(a) $41\,321 \div 844 \approx 49$ times. (b) One gridline apart would mean ten
times, two gridlines a hundred times; 49 is between, so the lines are more than
one gridline apart and less than two.
\end{solution}

\begin{solution}{exo:g6:living-standards:8}
(a) For example, a chart showing that another region --- East Asia, say ---
had an \omterm{def:g6:living-standards:income-per-person}{income per person} that rose steadily, decade after decade, from well
before 1800, while north-western Europe's stayed flat until later. (b) The rise
did not last: in 1800 the Netherlands had about the same \omterm{def:g6:living-standards:income-per-person}{income per person} as in
1600 ($4\,184$ against $4\,270$ dollars). A \omterm{def:g6:living-standards:take-off}{take-off} is a rise that goes on for
many decades.
\end{solution}

\begin{solution}{exo:g6:living-standards:9}
(a) $\dfrac{3\,000 \times 30\,000 + 7\,000 \times 10\,000}{10\,000}
= \dfrac{90\,000\,000 + 70\,000\,000}{10\,000} = 16\,000$ dalens.
(b) Nobody: 3\,000 people have $30\,000$ dalens and 7\,000 have $10\,000$.
(c) That the country is split in two: seven people in ten have much less than
the average, three in ten much more.
\end{solution}

\begin{solution}{exo:g6:living-standards:10}
(a) $1\,139.28$; $1\,150.67$; $1\,162.18$; $1\,173.80$; $1\,185.54$; $1\,197.40$;
$1\,209.37$; $1\,221.46$; $1\,233.67$; $1\,246.01$ dollars in 1830.
(b) $(1\,246.01 - 1\,128) \div 1\,128 = 118.01 \div 1\,128 \approx 0.105$: a rise
of about $10.5\,\%$. It is more than $10\,\%$ because each year's $1\,\%$ is taken
of a value that has already risen. (c) Each ten years multiplies by about
$1.105$: $1\,128 \times 1.105 \approx 1\,246.4$, then $1\,246.4 \times 1.105
\approx 1\,377.3$, then $1\,377.3 \times 1.105 \approx 1\,521.9$ dollars in 1850.
The source gives only $1\,301$ dollars: the world grew more slowly than $1\,\%$
a year.
\end{solution}

\begin{solution}{exo:g6:living-standards:11}
(a) $(18\,486 - 5\,267) \div 18\,486 = 13\,219 \div 18\,486 \approx 0.715$: a fall of
about $71.5\,\%$. (b) In 2022 it was below its 1950 level ($5\,267$ against
$8\,462$ dollars). (c) It is false: a country can lose, in ten years, more than
it gained in sixty.
\end{solution}

\begin{solution}{exo:g6:living-standards:12}
(a) Yes. Invented example: 100 people; at first each has 1\,000 dalens, so
\omterm{def:g6:living-standards:income-per-person}{income per person} is 1\,000. Thirty years later, 10 people have 11\,000 dalens
each and the other 90 still have 1\,000: \omterm{def:g6:living-standards:income-per-person}{income per person} is
$(110\,000 + 90\,000) \div 100 = 2\,000$, double, yet 90 people in 100 are no
better off. (b) Numbers on how the income is shared: for example the income of
the person in the middle, or of the poorest half, year after year. If \omterm{def:g6:living-standards:income-per-person}{income per
person} rises while the income of the middle person stays flat, the second half of
the sentence is true. One would also look at health and schooling, which \omterm{def:g6:living-standards:income-per-person}{income
per person} does not show.
\end{solution}

\begin{solution}{pb:g6:living-standards:1}
\textbf{1.} $12 \div 60 = 0.2$ dalen per hour.
\textbf{2.} In 1820: $0.4 \div 0.2 = 2$ hours. Today: $3 \div 18 = \frac{1}{6}$ hour,
that is 10 minutes. $120 \div 10 = 12$ times less.
\textbf{3.} $0.1 \div 5 = 0.02$ dalen per hour of candlelight; $0.02 \div 0.2 = 0.1$
hour of work, that is 6 minutes.
\textbf{4.} $0.02 \div 18 \approx 0.00111$ hour, and $0.00111 \times 3\,600 \approx 4$
seconds. $360 \div 4 = 90$ times less, for a brighter light.
\textbf{5.} 1950: Venezuela ($8\,462$), Argentina ($7\,949$), South Korea ($998$),
DR Congo ($909$), China ($799$), Botswana ($489$). 2022: South Korea ($41\,321$),
China ($19\,238$), Argentina ($18\,292$), Botswana ($16\,352$), Venezuela
($5\,267$), DR Congo ($844$).
\textbf{6.} South Korea: $41\,321 \div 998 \approx 41.4$ times. Argentina:
$18\,292 \div 7\,949 \approx 2.3$ times.
\textbf{7.} Venezuela: $(8\,462 - 5\,267) \div 8\,462 = 3\,195 \div 8\,462 \approx
0.378$, a fall of $37.8\,\%$. DR Congo: $(909 - 844) \div 909 = 65 \div 909 \approx
0.072$, a fall of $7.2\,\%$.
\textbf{8.} South Korea and China, which passed Argentina, and Botswana: in 1950
Argentina had $7\,949 \div 489 \approx 16.3$ times Botswana's \omterm{def:g6:living-standards:income-per-person}{income per person},
in 2022 only $18\,292 \div 16\,352 \approx 1.1$ times. Venezuela and DR Congo fell
further behind.
\textbf{9.} Fastland: 1\,000 (2000), 2\,000 (2010), 4\,000 (2020), 8\,000 (2030),
16\,000 (2040), 32\,000 (2050) dalens.
\textbf{10.} Slowland: 8\,000; $8\,000 \times 1.25 = 10\,000$; $12\,500$; $15\,625$;
$15\,625 \times 1.25 = 19\,531.25 \approx 19\,531$; $19\,531 \times 1.25 =
24\,413.75 \approx 24\,414$ dalens.
\textbf{11.} In 2040 Fastland is still behind ($16\,000 < 19\,531$); in 2050 it is
ahead ($32\,000 > 24\,414$). It passes Slowland between 2040 and 2050.
\textbf{12.} Two straight rising lines, because each country's income is
multiplied by the same number every ten years; Fastland's is steeper (it doubles,
Slowland's is multiplied by $1.25$). The lines of the six-paths chart are not
straight: they speed up, pause and sometimes fall. A race with fixed growth rates
is a simplification; real growth can stop or reverse, as Argentina's and Venezuela's
paths show.
\end{solution}
